% LuaLaTeX document — compile with: lualatex cell_averaged_coupling.tex  (twice)
\documentclass[11pt,a4paper]{article}

\usepackage{fontspec}
\setmainfont{Latin Modern Roman}
\setmonofont{Latin Modern Mono}

\usepackage{amsmath,amssymb,amsthm,bm}
\usepackage{booktabs}
\usepackage{geometry}
\geometry{margin=2.5cm}
\usepackage{hyperref}
\hypersetup{colorlinks=true,linkcolor=blue!60!black,citecolor=green!50!black,urlcolor=blue!70!black}
\usepackage{enumitem}
\usepackage{xcolor}
\usepackage{mathtools}

\newcommand{\D}{\Delta}
\newcommand{\order}{\mathcal{O}}
\newcommand{\dd}{\mathrm{d}}
\newcommand{\avg}[1]{\langle #1 \rangle}
\newcommand{\kpar}{k_{\parallel}}

\title{Cell-averaged coupling for a cubic-voxel lattice:\\
       why the midpoint rule cannot be repaired by summing further}
\author{}
\date{}

\begin{document}
\maketitle

\begin{abstract}
A voxel lattice couples its cells through a propagator that is almost always
\emph{sampled at cell centres}. The cell has volume, so this is a midpoint rule,
and its error is scale-invariant: refining the mesh shrinks the cell and the
spacing together, so the bias never dilutes. This note establishes that the
obvious repair --- correcting neighbours one at a time out to some radius ---
\emph{cannot be made to work}, because the correction has an $\order(1/R)$ tail
and becomes summation-shape dependent beyond $kr\sim1$, and the two requirements
are incompatible at any box size. It then gives a closed-form replacement in two
halves, exploiting the fact that the cube average is an identity on plane waves
but not on a spatial sum. Against an exact reference the corrected scheme
converges at second order where the truncated one saturates. A secondary
consequence is the closure of a long-standing empirical constant: a fitted shear
correction that had stood at $0.886234$, then $1.023737$, is measured here at
$1.0000$ and is shown to have been a propagator artefact throughout.
\end{abstract}

\tableofcontents

%======================================================================
\section{The defect}
%======================================================================

Write the coupling between two cells as the propagator $G$ evaluated at the
separation of their centres. The source cell carries a volume moment, so the
field its neighbour should receive is
%
\begin{equation}
  \psi_i \;=\; \Bigl[\frac{1}{V}\int_{V_{\text{source}}} G(\bm{x}_i - \bm{x}')\,\dd^3x'\Bigr]
              \times (\text{moment}),
  \label{eq:source-average}
\end{equation}
%
the average over the source cell, evaluated at the \emph{receiver centre}. The
usual practice replaces the bracket by $G(\bm{x}_i-\bm{x}_j)$: a midpoint rule.

The error is
%
\begin{equation}
  \avg{G} - G \;=\; \frac{d^{2}}{24}\,\nabla^{2}G \;+\; \order(d^{4}),
  \label{eq:midpoint-error}
\end{equation}
%
with $d$ the cell side. What makes this pernicious rather than merely
inaccurate is that it is \emph{scale-invariant on a space-filling lattice}. The
cells and their separations shrink together, so the relative error at each
neighbour is unchanged by refinement, and the summed bias tends to a constant
that no mesh refinement removes.

\paragraph{Which average.}
Equation~\eqref{eq:source-average} averages over the source cell only, with the
receiver a point. That is the choice consistent with a collocation scheme, whose
state variable is a point value at the cell centre. A Galerkin scheme, whose
state is a cell average, wants a \emph{double} average over both cells. The two
differ, and the distinction is load-bearing: on this solver the collocation
pairing outperforms the Galerkin one by factors of $3.4$ and $16$ against an
exact reference. For an ellipsoid the two coincide, by Eshelby's uniformity
theorem; a cube is the only place they can differ at all.

%======================================================================
\section{Why the obvious repair fails}
%======================================================================

The natural fix is to correct neighbours directly out to some radius $R$ and
accept the midpoint rule beyond it. This fails for two reasons that are
individually survivable and jointly fatal.

\subsection{The tail is long-ranged}

Away from the origin the propagator entries satisfy the Helmholtz equation,
$\nabla^{2}g_c = -k_c^{2}g_c$ exactly, and derivatives commute with
$\nabla^{2}$. Two regimes follow from \eqref{eq:midpoint-error}:

\begin{itemize}[leftmargin=2em]
\item $kr \ll 1$: the elastostatic Kelvin tensor is \emph{biharmonic}, not
  harmonic ($\nabla^{2}r = 2/r$), so $\nabla^{2}G \neq 0$ but sits two orders
  below $G$. With the displacement block $G\sim1/r$ this gives a correction
  $\sim d^{2}/r^{3}$, a shell sum $\sim d^{2}/R^{2}$, and an $\order(1/R)$
  tail --- slow, but absolutely convergent.
\item $kr \gg 1$: Helmholtz makes the correction $\to -(kd)^{2}G/24$, so it
  \emph{tracks the kernel itself}. Each shell then contributes a constant and
  only oscillation converges the sum --- conditionally.
\end{itemize}

Both halves are measured. The far-field ratio $|\avg{G}-G|/|G|$ plateaus at
$9.14\times10^{-4}$ against the predicted $(k_Sd)^{2}/24 = 9.375\times10^{-4}$,
within $3\%$; and the static ratio falls as $1/r^{2}$ exactly, the signature of a
surviving $d^{2}\nabla^{2}$ term rather than a cancelled one.

\subsection{The far field is shape-dependent}

Summing the correction over a \emph{square} region and over a \emph{circular}
one tests conditional convergence directly. At $k_S r_{\max}=0.4$ the gap
between them shrinks by $5.4\times$: absolutely convergent. At
$k_S r_{\max}=3$ it \emph{grows} by $2\times$: conditionally convergent, so the
limit depends on the shape of the region summed.

\subsection{The two are incompatible}

\begin{quote}
Converging the $\order(1/R)$ tail requires $R\sim80$ cells. The shape term
switches on near $R \sim 1/(k_S d) \approx 50$ cells. \emph{The box cannot be
made large enough to converge without becoming large enough to be
shape-dependent.}
\end{quote}

No radius is correct. This is a structural obstruction, not a tuning problem,
and it is why the truncated correction had to be replaced rather than extended.

%======================================================================
\section{The closed-form replacement}
%======================================================================

The cell average is a convolution with the cell indicator, so on a plane wave it
is exactly a multiplication by the cube form factor
%
\begin{equation}
  f_0(\bm{k}) \;=\; \operatorname{sinc}(k_1h)\,\operatorname{sinc}(k_2h)\,
                    \operatorname{sinc}(k_3h), \qquad h=d/2,
  \label{eq:form-factor}
\end{equation}
%
\emph{one} power --- the single average of \eqref{eq:source-average}. The
squared form factor $|f_0|^{2}$, which appears in the existing literature and in
this package's own inter-voxel tables, is the \emph{double} (Galerkin) average:
source cell convolved with field cell.

\paragraph{A form factor multiplies plane waves, and nothing else.}
This is the organising fact, and it splits the construction in two. Where the
lattice kernel is a plane-wave sum, the average is an identity. Where it is a
spatial summation, no form factor exists.

\subsection{Between planes: an identity}

For non-equal depths the Bloch kernel is a plane-wave sum, and
\eqref{eq:form-factor} applies directly. Two properties of the implementation
are worth recording because neither is optional.

\emph{Per mode.} The $P$ and $S$ poles carry different $k_3$, so a single scalar
factor on the assembled block merely rescales it; each pole must be scaled by
the sinc built from its own $k_3$.

\emph{The decay margin halves.} On the evanescent branch $k_3 = i\kappa$ and
%
\begin{equation}
  \operatorname{sinc}(i\kappa h) \;=\; \frac{\sinh \kappa h}{\kappa h}
  \;\sim\; \frac{e^{\kappa h}}{2\kappa h},
\end{equation}
%
which \emph{grows}. Against the kernel's own $e^{-\kappa|\D z|}$ the product
decays as $e^{-\kappa(|\D z|-h)}$: still convergent, because $|\D z|\ge 2h$ for
distinct planes, but at half the rate at adjacent planes. A reciprocal-sum
cutoff tuned for the point kernel silently returns $10^{-8}$ where it had
returned $10^{-16}$.

\subsection{In-plane: an analytic tail}

At equal depth the route is an Ewald sum, whose real-space half is a spatial
summation --- so no form factor applies. Nor can the plain reciprocal sum be
substituted: with the \emph{single} form factor its terms decay as $1/|\bm{G}|^{2}$
against $\sim|\bm{G}|$ growth of the strain block and $\sim|\bm{G}|$ states per
shell, which is logarithmically divergent. (The double form factor gains a
further $1/|\bm{G}|^{2}$ and does converge, which is why the squared route never
met this obstruction.)

Instead, using $\nabla^{2}g_c = -k_c^{2}g_c$ exactly away from the origin,
%
\begin{equation}
  \sum_{|\bm{R}|>R_0} \avg{D} \;=\;
  \Bigl(1 - \frac{k_c^{2}d^{2}}{24}\Bigr)\bigl(D_{\text{full}} - D_{\text{near}}\bigr)
  \;+\; \order(d^{4}),
  \label{eq:tail}
\end{equation}
%
where $D_{\text{full}}$ is the existing Ewald sum --- already computed once per
mode, which is exactly the granularity \eqref{eq:tail} needs, since $P$ and $S$
carry different $k^{2}$. The near shell is averaged directly; the far field is
closed form.

\paragraph{The sharp test is $R_0$-independence.} $R_0$ splits the sum between
the two pieces, so a wrong tail would make the answer drift with it. Measured
drift is $3.2\times10^{-6}$ and $3.8\times10^{-6}$ for $R_0=3,4$ against $R_0=2$.

%======================================================================
\section{Results}
%======================================================================

Against Kennett, exact for a uniform layer and independent of the construction:

\begin{table}[h]
\centering
\begin{tabular}{llrr}
\toprule
case & scheme & \multicolumn{2}{c}{error ratio per mesh halving}\\
\cmidrule(lr){3-4}
 & & $n_z\!:\!2\to4$ & $4\to8$ \\
\midrule
pure shear & truncated shell & $2.2$ & $1.29$\\
           & closed form     & $6.1$ & $3.70$\\
mixed      & truncated shell & $2.0$ & $1.19$\\
\bottomrule
\end{tabular}
\caption{The truncated scheme \emph{saturates} --- its ratio tends to $1.0$,
a floor --- because its bias is scale-invariant. The closed-form scheme tends
to $4.0$, second order.}
\end{table}

The single-mesh error is the least informative number here; the distinction is
between a scheme that keeps improving under refinement and one that stops. The
mixed case's closed-form column is omitted deliberately: its error passes
through zero between $n_z=4$ and $n_z=8$, so no convergence order can be quoted
for it. Its magnitude is below the truncated scheme's at every mesh tested.

\paragraph{Validity.} The in-plane tail is an expansion in $(\kappa d)^{2}$.
Measured across frequency it is flat to $k_Sd\approx1$ and degrades beyond;
a larger near shell reduces the error at every frequency tested, confirming that
the expansion and not something else is the limiting term. The regime of use,
$ka<0.3$, sits an order of magnitude inside that.

%======================================================================
\section{A fitted constant, closed}
%======================================================================

A one-parameter shear correction $K$ had been carried for some time, defined so
that $K=1$ means the cube's own shear response is correct. Its history:
%
\begin{center}
\begin{tabular}{lll}
\toprule
value & context & what it was measuring\\
\midrule
$0.886234$ & original & a doubly-averaged contact term against a collocation state\\
$1.023737$ & contact fixed & the remaining average truncated at a finite shell\\
$1.0000$   & this work & nothing\\
\bottomrule
\end{tabular}
\end{center}
%
Measured on the corrected propagator, $|K-1| = 2.75\times10^{-4}$. Each time a
propagator defect was removed, $K$ moved toward unity; there is no mechanism to
derive. The near-misses $8/9 = 0.888889$ and $\sqrt{\pi}/2 = 0.886227$ were
close to an artefact.

\paragraph{The methodological point is the durable one.} A one-parameter
correction fitted against the same reference it is then scored on will look
successful: $K$ was universal across contrasts, did nothing to the channels it
should not have touched, and improved cases it was never fitted to. It was still
an artefact. Relatedly, the tempting success metric --- how much the answer
improves once the parameter is tuned --- reaches $290\times$ here and means
nothing, because one knob with leverage cancels almost any small scalar
residual. The honest measure is the distance of the \emph{unfitted} prediction
from ideal.

%======================================================================
\section{What is not established}
%======================================================================

\begin{enumerate}[leftmargin=2em]
\item \textbf{The in-plane half is still an expansion.} Equation~\eqref{eq:tail}
  is exact through $d^{2}$; the between-plane half is exact at any cell size.
  This is the one place a small-$d$ limit re-enters a formulation that otherwise
  avoids them, and it is the origin of the frequency degradation above.
  Section~\ref{sec:exact-inplane} sets out what an exact version would require.
\item \textbf{The dynamic verification is incomplete.} The spectral single
  average was verified against a real-space reference in the static limit, where
  the ultraviolet behaviour in question lives. The dynamic case, with its poles
  and the subtracted formulation they force, is untested.
\item \textbf{One convergence order is unpinned}, as noted above.
\item \textbf{The isolated cube is not the lattice.} Against a subdivided-cube
  reference that resolves the internal field, the uniform-strain single-site
  T-matrix differs by $\approx0.35\%$, while its effective response in a
  space-filling lattice is correct to $2.75\times10^{-4}$. Both hold: the
  discrepancy largely cancels among identical neighbours. The $0.35\%$ is
  \emph{not} decisive --- changing the sub-cell coupling moves it by four times
  its own size --- and sharpening it requires a coupling whose bias is below
  $10^{-3}$, not a richer single-site basis.
\end{enumerate}

%======================================================================
\section{Making the in-plane average exact}
\label{sec:exact-inplane}
%======================================================================

The obstruction is stated precisely in \eqref{eq:tail}: the Ewald real-space
half is a spatial sum, so the form factor does not apply, and the plain
reciprocal sum is logarithmically divergent under the single form factor. The
$d^{2}$ expansion is a way around that, not a resolution of it.

Averaging is linear, so it distributes across the Ewald split:
%
\begin{equation}
  \sum_{\bm{R}} \avg{G}(\bm{R})\,e^{-i\bm{k}\cdot\bm{R}}
  \;=\; \avg{\text{real-space screened half}} \;+\; \avg{\text{reciprocal half}} .
\end{equation}
%
Each half admits an exact treatment, by different means:

\begin{itemize}[leftmargin=2em]
\item \textbf{Reciprocal half.} Plane waves, so the form factor multiplies. The
  vertical factor is not $\operatorname{sinc}$ at equal depth, because the
  average crosses the $|z|$ branch point at the cell centre, but it remains in
  closed form:
  \begin{equation}
    \frac{1}{2h}\int_{-h}^{h} e^{i\gamma|u|}\,\dd u
      \;=\; \frac{e^{i\gamma h}-1}{i\gamma h}.
  \end{equation}
\item \textbf{Real-space half.} The screened kernel is short-ranged and regular
  at every non-self lattice site --- the nearest in-plane neighbour has
  $|\bm r| \ge d - h\sqrt{3} > 0$, and $\bm R = 0$ is excluded as the self term
  --- so ordinary quadrature averages it exactly. No form factor is needed, and
  none exists.
\end{itemize}

This would remove the last small-$d$ expansion from the construction and with it
the frequency degradation, at the cost of reworking the Ewald internals rather
than wrapping them. It is not implemented here, and the argument above is a
derivation sketch, not a validated result.

\end{document}
